%========================================
% MatinBook — Stage 11: Cross-References Test
% Version: v1.1
% Purpose:
%   Validate the cross-reference system of MatinBook v1.1:
%     1. \ref (standard reference, no prefix)
%     2. \cref (cleveref, auto-detects the type)
%     3. \Cref (capitalized form, same as \cref in Persian)
%     4. \crefrange (range reference)
%     5. Multiple references (\cref{a,b,c})
%     6. Helper commands (\thmref, \lemref, \corref,
%        \propref, \defref, \exref, \remref, \excref)
%     7. Cross-chapter references
%     8. Hyperref bookmarks and links
%
% Compilation:
%   xelatex -shell-escape stage11-references.tex
%   xelatex -shell-escape stage11-references.tex
%   xelatex -shell-escape stage11-references.tex
%   (3 passes recommended for stable cross-references)
%
% Expected outcome:
%   A professionally typeset PDF demonstrating that every
%   reference resolves to the correct number and Persian
%   prefix, with working hyperlinks and no "??" markers.
%========================================

%------------------------------------------------------------------------
% Search Path (for tests directory)
%------------------------------------------------------------------------
\makeatletter
\def\input@path{{../}}
\makeatother

%------------------------------------------------------------------------
% Document Class
%------------------------------------------------------------------------
\documentclass{matinbook}

%------------------------------------------------------------------------
% Book Metadata
%------------------------------------------------------------------------
\title{آزمون سیستم ارجاع هوشمند}
\author{تیم توسعه MatinBook}
\date{مهر ۱۴۰۵}

\booktitle{آزمون سیستم ارجاع هوشمند}

\renewcommand{\repository}{github.com/matinmahmoudi-cs/matinbook}

%========================================================================
% DOCUMENT
%========================================================================
\begin{document}

%------------------------------------------------------------------------
% FRONT COVER
%------------------------------------------------------------------------
\makecover
    {آزمون سیستم ارجاع هوشمند}
    {ارزیابی \lr{hyperref} و \lr{cleveref} در MatinBook}
    {تیم توسعه MatinBook}
    {مهر ۱۴۰۵}

%------------------------------------------------------------------------
% FRONT MATTER (symmetric margins)
%------------------------------------------------------------------------
\frontmatter
\frontmattergeometry

% Title page
\maketitle

% Copyright / Colophon
\thispagestyle{empty}
\vfill
\begin{center}
    {\large\textbf{شناسنامه آزمون}}

    \vspace{0.8cm}

    \begin{tabular}{rl}
        عنوان: & آزمون سیستم ارجاع هوشمند \\
        نسخه: & \matinbookversion \\
        تاریخ: & \matinbookdate \\
        موتور: & \lr{XeLaTeX} \\
        مجوز: & \lr{MIT} \\
    \end{tabular}

    \vspace{0.8cm}

    {\small این سند بخشی از مجموعه آزمون‌های یکپارچگی \lr{MatinBook} است.}
\end{center}
\clearpage

% Preface
\chapter*{پیشگفتار}
\addcontentsline{toc}{chapter}{پیشگفتار}

این سند، یازدهمین آزمون از مجموعه‌ی آزمون‌های یکپارچگی
\lr{MatinBook} نسخه‌ی \lr{\matinbookversion} است. هدف آن،
اعتبارسنجی کامل سیستم ارجاع هوشمند کلاس در هشت حوزه‌ی
متمایز است: \lr{\textbackslash ref} استاندارد،
\lr{\textbackslash cref}، \lr{\textbackslash Cref}،
\lr{\textbackslash crefrange}، ارجاع‌های چندگانه،
دستورات کمکی، ارجاع‌های بین‌فصلی و بوکمارک‌های
\lr{hyperref}.

در کتاب‌های فنی و علمی، ارجاع دقیق به عناصر مختلف
(معادله، قضیه، شکل، جدول، بخش) اهمیت ویژه‌ای دارد.
خواننده باید بتواند به‌سرعت به مرجع مورد نظر برسد و
از پیوستگی مطالب مطمئن شود. \lr{MatinBook} با
استفاده از \lr{cleveref} و دستورات سفارشی، این
قابلیت‌ها را فراهم می‌سازد.

هر بخش از این آزمون، با مثال‌های عملی و ارجاعات متقابل
همراه است تا خواننده بتواند درستی رفتار هر زیرسیستم را
در بافت واقعی یک کتاب ارزیابی کند.

\clearpage

% Table of Contents
\tableofcontents
\clearpage

% List of Figures
\listoffigures
\clearpage

% List of Tables
\listoftables
\clearpage

%------------------------------------------------------------------------
% MAIN MATTER (wide margin for notes)
%------------------------------------------------------------------------
\mainmatter
\mainmattergeometry

%========================================================================
% CHAPTER 1 — CREATING REFERENCE TARGETS
%========================================================================
\chapter{ایجاد اهداف ارجاع}
\label{chap:targets}

\section{معادلات}
\label{sec:targets-equations}

چند معادله‌ی معروف که به آن‌ها ارجاع خواهیم داد:

\begin{equation}
    E = mc^{2}
    \label{eq:einstein}
\end{equation}

\begin{equation}
    a^{2} + b^{2} = c^{2}
    \label{eq:pythagoras}
\end{equation}

\begin{equation}
    e^{i\pi} + 1 = 0
    \label{eq:euler}
\end{equation}

\mnote{معادلات با فرمت «شماره-فصل» شماره‌گذاری می‌شوند، مطابق استاندارد نشر ایران.}

\section{قضایا و تعاریف}
\label{sec:targets-theorems}

\begin{theorem}[قضیه‌ی مقدار میانی]
    اگر تابع $f$ روی بازه‌ی بسته‌ی $[a, b]$ پیوسته باشد
    و $y$ عددی بین $f(a)$ و $f(b)$ باشد، آنگاه حداقل یک
    $c \in [a, b]$ وجود دارد که $f(c) = y$.
    \label{thm:intermediate-value}
\end{theorem}

\begin{definition}[پیوستگی]
    تابع $f$ در نقطه‌ی $x_{0}$ \textbf{پیوسته} است اگر:
    \[
        \lim_{x \to x_{0}} f(x) = f(x_{0})
    \]
    \label{def:continuity}
\end{definition}

\begin{lemma}[لم فشردگی]
    هر دنباله‌ی کراندار در $\R$ دارای زیردنباله‌ی همگراست.
    \label{lem:bolzano-weierstrass}
\end{lemma}

\begin{corollary}
    هر دنباله‌ی یکنوا و کراندار همگراست.
    \label{cor:monotone-convergence}
\end{corollary}

\begin{proposition}
    مجموعه‌ی اعداد گویا $\Q$ در $\R$ چگال است.
    \label{prop:q-dense}
\end{proposition}

\section{مثال‌ها و نکات}
\label{sec:targets-examples}

\begin{example}[تابع پیوسته]
    تابع $f(x) = x^{2}$ روی تمام $\R$ پیوسته است،
    زیرا برای هر $x_{0}$:
    \[
        \lim_{x \to x_{0}} x^{2} = x_{0}^{2} = f(x_{0})
    \]
    \label{ex:continuous-function}
\end{example}

\begin{remark}
    عکس \cref{thm:intermediate-value} لزوماً برقرار نیست.
    یعنی ممکن است تابعی در شرط مقدار میانی صدق کند اما
    پیوسته نباشد.
    \label{rem:ivt-converse}
\end{remark}

\section{شکل و جدول}
\label{sec:targets-figure-table}

نمودار \cref{fig:parabola} تابع درجه دو را نشان می‌دهد:

\begin{figure}[h]
\centering
\begin{latin}
\begin{tikzpicture}
  \begin{axis}[
    width=\textwidth, height=6cm,
    xlabel={$x$},
    ylabel={$f(x)$},
    title={Continuous Function},
    grid=major,
    domain=-2:2,
    samples=100
  ]
    \addplot[matinblue, thick] {x^2};
    \addplot[matinred, only marks, mark=*] coordinates {(1, 1)};
    \node[above] at (axis cs:1,1) {$(1, f(1))$};
  \end{axis}
\end{tikzpicture}
\end{latin}
\caption{نمودار تابع $f(x) = x^{2}$}
\label{fig:parabola}
\end{figure}

جدول \cref{tab:parabola-values} مقادیر عددی این تابع را
فهرست می‌کند:

\begin{matintable}{مقادیر تابع $f(x) = x^{2}$}{tab:parabola-values}
\begin{tblr}{
    width = 0.5\textwidth,
    colspec = {c c},
    hlines, vlines,
    colsep = 10pt,
    rowsep = 4pt,
    row{1} = {font=\bfseries},
}
$x$ & $f(x)$ \\
$-2$ & $4$ \\
$-1$ & $1$ \\
$0$ & $0$ \\
$1$ & $1$ \\
$2$ & $4$ \\
\end{tblr}
\end{matintable}

%========================================================================
% CHAPTER 2 — TESTING ALL REFERENCE COMMANDS
%========================================================================
\chapter{آزمون تمام دستورات ارجاع}
\label{chap:references}

\section{ارجاع استاندارد با \lr{\textbackslash ref}}
\label{sec:ref-standard}

دستور \lr{\textbackslash ref} فقط شماره‌ی عنصر را
بدون پیشوند نمایش می‌دهد:

\begin{itemize}
    \item معادله‌ی اینشتین: \ref{eq:einstein}
    \item قضیه‌ی مقدار میانی: \ref{thm:intermediate-value}
    \item نمودار سهمی: \ref{fig:parabola}
\end{itemize}

\section{ارجاع هوشمند با \lr{\textbackslash cref}}
\label{sec:ref-cleveref}

دستور \lr{\textbackslash cref} به‌صورت خودکار نوع عنصر
را تشخیص می‌دهد و پیشوند مناسب اضافه می‌کند:

\begin{itemize}
    \item \cref{eq:einstein} معادله‌ی معروف اینشتین است.
    \item طبق \cref{thm:intermediate-value}، تابع پیوسته
    شرط مقدار میانی را دارد.
    \item در \cref{fig:parabola} نمودار سهمی را می‌بینیم.
    \item \cref{tab:parabola-values} مقادیر تابع را نشان
    می‌دهد.
    \item \cref{def:continuity} تعریف پیوستگی است.
    \item \cref{ex:continuous-function} یک مثال است.
    \item \cref{rem:ivt-converse} یک نکته‌ی مهم است.
    \item \cref{lem:bolzano-weierstrass} لم فشردگی است.
    \item \cref{cor:monotone-convergence} نتیجه‌ی لم است.
    \item \cref{prop:q-dense} گزاره‌ی چگالی $\Q$ است.
\end{itemize}

\section{ارجاع اختصاصی MatinBook}
\label{sec:ref-helpers}

این دستورات برای راحتی بیشتر تعریف شده‌اند و همه به
\lr{\textbackslash cref} واگذار می‌شوند:

\begin{itemize}
    \item \textbf{قضیه:} \cref{thm:intermediate-value}
    پایه‌ی آنالیز است.
    \item \textbf{تعریف:} \cref{def:continuity} مفهوم
    حد را دقیق می‌کند.
    \item \textbf{لم:} \cref{lem:bolzano-weierstrass}
    در اثبات‌ها کاربرد دارد.
    \item \textbf{نتیجه:} \cref{cor:monotone-convergence}
    از لم نتیجه می‌شود.
    \item \textbf{گزاره:} \cref{prop:q-dense} چگالی
    $\Q$ در $\R$ را بیان می‌کند.
    \item \textbf{مثال:} \cref{ex:continuous-function}
    ساده ولی گویاست.
    \item \textbf{نکته:} \cref{rem:ivt-converse} هشدار
    مهمی می‌دهد.
\end{itemize}

> **نکته:** دستورات \lr{\textbackslash meqref}،
> \lr{\textbackslash figref}، \lr{\textbackslash tabref}،
> \lr{\textbackslash chref} و \lr{\textbackslash secref}
> در v1.1 حذف شده‌اند. به‌جای آن‌ها از \lr{\textbackslash cref}
> استفاده کنید.

\section{ارجاع‌های هوشمند در عمل}
\label{sec:ref-in-practice}

حالا از ترکیب این ارجاع‌ها در یک متن واقعی استفاده می‌کنیم:

\medskip

طبق \cref{thm:intermediate-value}، اگر تابعی پیوسته باشد
(\cref{def:continuity})، آنگاه شرط مقدار میانی را دارد.
\cref{ex:continuous-function} نشان می‌دهد که $x^{2}$
(\cref{fig:parabola}) این شرط را دارد. اما
\cref{rem:ivt-converse} هشدار می‌دهد که عکس
\cref{thm:intermediate-value} برقرار نیست. برای مطالعه‌ی
بیشتر، به \cref{chap:cross-chapter} مراجعه کنید.

%========================================================================
% CHAPTER 3 — RANGE REFERENCES
%========================================================================
\chapter{ارجاع‌های بازه‌ای}
\label{chap:range-refs}

\section{ارجاع به چند معادله}
\label{sec:range-multiple}

قابلیت \lr{cleveref} امکان ارجاع به چند عنصر را
به‌صورت فشرده فراهم می‌کند:

\begin{itemize}
    \item \cref{eq:einstein,eq:pythagoras,eq:euler}
    سه معادله‌ی معروف هستند.
    \item طبق \cref{thm:intermediate-value,lem:bolzano-weierstrass,cor:monotone-convergence}
    نتایج مهمی در آنالیز داریم.
\end{itemize}

\section{ارجاع بازه‌ای}
\label{sec:range-range}

با استفاده از \lr{\textbackslash crefrange}:

\begin{itemize}
    \item معادلات \crefrange{eq:einstein}{eq:euler}
    همگی از معادلات بنیادی هستند.
\end{itemize}

%========================================================================
% CHAPTER 4 — CROSS-CHAPTER REFERENCES
%========================================================================
\chapter{ارجاع بین‌فصلی}
\label{chap:cross-chapter}

\section{ارجاع به فصل‌های دیگر}
\label{sec:cross-chapters}

حالا می‌خواهیم به عناصری که در \cref{chap:targets}
تعریف کردیم، از این فصل ارجاع دهیم:

\begin{itemize}
    \item معادله‌ی \cref{eq:pythagoras} که در
    \cref{chap:targets} تعریف شد، اساس هندسه است.
    \item \cref{thm:intermediate-value} (از
    \cref{chap:targets}) یکی از قضایای اساسی آنالیز
    ریاضی است.
    \item \cref{fig:parabola} در \cref{chap:targets}
    نمودار تابع درجه دو را نشان می‌دهد.
\end{itemize}

\section{ارجاع به این فصل از فصول قبلی}
\label{sec:cross-backward}

همچنین می‌توان از فصول قبلی به این فصل ارجاع داد.
برای مثال، اگر در \cref{chap:targets} بودیم،
می‌توانستیم بنویسیم:
«در \cref{chap:cross-chapter} ارجاع بین‌فصلی را
بررسی خواهیم کرد.»

%========================================================================
% CHAPTER 5 — SUMMARY
%========================================================================
\chapter{خلاصه‌ی نتایج}
\label{chap:summary}

\section{وضعیت سیستم ارجاع}
\label{sec:summary-status}

جدول \cref{tab:reference-status} وضعیت تمام دستورات
ارجاع \lr{MatinBook} را نشان می‌دهد.

\begin{matintable}{وضعیت تمام دستورات ارجاع}{tab:reference-status}
\begin{tblr}{
    width = \textwidth,
    colspec = {l X[c] X[c]},
    hlines, vlines,
    colsep = 8pt,
    rowsep = 4pt,
    row{1} = {font=\bfseries},
}
دستور & کاربرد & وضعیت \\
\lr{\textbackslash ref} & شماره‌ی عنصر بدون پیشوند & \textcolor{matingreen}{فعال} \\
\lr{\textbackslash cref} & شماره با پیشوند هوشمند & \textcolor{matingreen}{فعال} \\
\lr{\textbackslash Cref} & مشابه \lr{\textbackslash cref} (فارسی) & \textcolor{matingreen}{فعال} \\
\lr{\textbackslash crefrange} & ارجاع بازه‌ای & \textcolor{matingreen}{فعال} \\
\lr{\textbackslash cref\{a,b,c\}} & ارجاع چندگانه & \textcolor{matingreen}{فعال} \\
\lr{\textbackslash thmref} & ارجاع به قضیه & \textcolor{matingreen}{فعال} \\
\lr{\textbackslash lemref} & ارجاع به لم & \textcolor{matingreen}{فعال} \\
\lr{\textbackslash corref} & ارجاع به نتیجه & \textcolor{matingreen}{فعال} \\
\lr{\textbackslash propref} & ارجاع به گزاره & \textcolor{matingreen}{فعال} \\
\lr{\textbackslash defref} & ارجاع به تعریف & \textcolor{matingreen}{فعال} \\
\lr{\textbackslash exref} & ارجاع به مثال & \textcolor{matingreen}{فعال} \\
\lr{\textbackslash remref} & ارجاع به نکته & \textcolor{matingreen}{فعال} \\
\lr{\textbackslash excref} & ارجاع به تمرین & \textcolor{matingreen}{فعال} \\
\end{tblr}
\end{matintable}

\section{معیارهای ارزیابی}
\label{sec:summary-criteria}

در این آزمون، سیستم ارجاع \lr{MatinBook} بر اساس پنج
معیار ارزیابی شد:

\begin{enumerate}
    \item \textbf{دقت}: شماره‌ی درست برای هر عنصر.
    \item \textbf{پیشوند فارسی}: «قضیه»، «شکل»، «معادله» و...
    \item \textbf{ارجاع بازه‌ای}: \lr{\textbackslash crefrange}.
    \item \textbf{ارجاع چندگانه}: \lr{\textbackslash cref\{a,b,c\}}.
    \item \textbf{لینک فعال}: \lr{hyperref} با رنگ آبی.
\end{enumerate}

\section{نتیجه‌گیری}
\label{sec:summary-conclusion}

\textbf{تمام زیرسیستم‌های ارجاع هوشمند \lr{MatinBook}
نسخه‌ی \lr{\matinbookversion} با موفقیت بارگذاری و
آزمایش شدند. سیزده دستور ارجاع شامل \lr{hyperref}،
\lr{cleveref} و دستورات کمکی، همگی به‌درستی کار می‌کنند
و شماره‌ها با پیشوندهای فارسی نمایش داده می‌شوند.}

%------------------------------------------------------------------------
% BACK MATTER (symmetric margins)
%------------------------------------------------------------------------
\backmatter
\frontmattergeometry

% Bibliography (empty placeholder)
\printbibliography[title={منابع و مراجع}]

% Index
\printindex

%------------------------------------------------------------------------
% BACK COVER
%------------------------------------------------------------------------
\authorbio{%
    تیم توسعه‌ی \lr{MatinBook}، گروهی از پژوهشگران و برنامه‌نویسان
    فارسی‌زبان است که با هدف ارائه‌ی یک راه‌حل حرفه‌ای برای
    حروف‌چینی کتاب‌های فنی فارسی فعالیت می‌کند. این پروژه حاصل
    سال‌ها تجربه در حوزه‌ی نشر دیجیتال و برنامه‌نویسی است.
}

\makebackcover{%
    این سند، یازدهمین آزمون از مجموعه‌ی آزمون‌های یکپارچگی
    \lr{MatinBook} نسخه‌ی \lr{\matinbookversion} است.

    \vspace{0.5cm}

    \textbf{زیرسیستم‌های آزموده‌شده:}
    \begin{itemize}
        \item \lr{\textbackslash ref} استاندارد
        \item \lr{\textbackslash cref} هوشمند
        \item \lr{\textbackslash crefrange} بازه‌ای
        \item ارجاع‌های چندگانه
        \item دستورات کمکی
        \item ارجاع‌های بین‌فصلی
        \item لینک‌های \lr{hyperref}
    \end{itemize}
}

\end{document}